\renewcommand{\Anno}{1999}

\renewcommand{\NumeroProblema}{3}

\renewcommand{\Testo}{Un tizio gioca nel seguente modo.  Punta su un
numero compreso tra 1 e 6, poi tira un dado perfettamente equilibrato. Supponiamo che esca il punteggio $n$.  Egli, allora, ritira il dado fino a quando non esce un punteggio $k\leq n$.  Se il numero $k$ è quello su cui aveva puntato, il giocatore vince e guadagna $k$ gettoni, altrimenti non vince e non guadagna nulla.
	
Determinare, a lungo andare, su quale numero conviene puntare.}

\renewcommand{\Risposta}{Conviene puntare sul numero 2.}

\renewcommand{\Aiutino}{Occorre calcolare, per ogni numero tra 1 e 6,
il guadagno medio del giocatore.  Il guadagno medio si calcola come il prodotto tra la probabilità di vittoria ed il guadagno in caso di vittoria.  Per calcolare la probabilità di vittoria dopo aver puntato su un certo numero, occorre distinguere i vari casi.  Ad esempio, puntando sul 4, si vince se $n=4$ (probabilità 1/6) e poi $k=4$ (probabilità 1/4), oppure se $n=5$ (probabilità 1/6) e poi $k=4$ (probabilità 1/5), oppure se $n=6$ (probabilità 1/6)
e poi $k=4$ (probabilità 1/6).}

\renewcommand{\SoluzioneUfficiale}{
}

\renewcommand{\SoluzioniAlternative}{
}

\renewcommand{\Osservazioni}{
}

\renewcommand{\Commenti}{
}
% Aggiunta redazionale: la soluzione preesistente non viene modificata.
\providecommand{\SoluzioneRedazionale}{}
\renewcommand{\SoluzioneRedazionale}{
% SSSUP-BOZZA-AI: SoluzioneRedazionale; autore=ChatGPT; data=2026-11-08; stato=DA-VERIFICARE
% SSSUP-PROVENIENZA: AI: elaborazione nuova dal testo
Se il primo lancio vale $n$, il risultato del primo lancio successivo non superiore a $n$ è uniforme su $\{1,\ldots,n\}$. Infatti, per ciascun $j\leq n$, la probabilità che sia proprio $j$ è
\[
\sum_{r=0}^{\infty}\left(\frac{6-n}{6}\right)^r\frac16=\frac1n.
\]
Puntando sul numero $j$, il guadagno medio per partita è pertanto
\[
G(j)=\frac j6\sum_{n=j}^6\frac1n.
\]
Per $1\leq j\leq5$,
\[
G(j+1)-G(j)=\frac16\left(\sum_{n=j+1}^6\frac1n-1\right).
\]
Per $j=1$ la somma è $29/20>1$, mentre per $j=2$ è $19/20<1$ e per indici maggiori diminuisce ulteriormente. Dunque $G$ cresce da $1$ a $2$ e poi decresce: conviene puntare sul $2$, con guadagno medio $G(2)=29/60$ gettoni per partita.
}
